\SourceLeaf{258}%
as follows. Differentiating the first and second of the equations given at the beginning of this article, we obtain
\[
\sin\psi\mathord{.}\frac{\mathrm{d}n}{\mathrm{d}q}+n\cos\psi\mathord{.}\frac{\mathrm{d}\psi}{\mathrm{d}q}+\sin\psi\mathord{.}\frac{\mathrm{d}\psi}{\mathrm{d}p}=0
\]
which, combined with
\[
n\cos\psi\mathord{.}\frac{\mathrm{d}\varphi}{\mathrm{d}q}+\sin\psi\mathord{.}\frac{\mathrm{d}\varphi}{\mathrm{d}p}=0
\]
gives
\[
\frac{r\sin\psi}{n}\mathord{.}\frac{\mathrm{d}n}{\mathrm{d}q}+\frac{r\sin\psi}{n}\mathord{.}\frac{\mathrm{d}(\psi+\varphi)}{\mathrm{d}p}+r\cos\psi\mathord{.}\frac{\mathrm{d}(\psi+\varphi)}{\mathrm{d}q}=0
\]
From this equation we can easily obtain the series for \(\psi+\varphi\) by the method of undetermined coefficients, bearing in mind that its first term must be \(\tfrac{1}{2}\pi\) when the radius is taken as the unit and \(2\pi\) denotes the circumference of the circle,
\[
\begin{aligned}
\relax[6]\quad\psi+\varphi={}&\tfrac{1}{2}\pi-f^0pq-\tfrac{2}{3}f'ppq-(\tfrac{1}{2}f''-\tfrac{1}{6}f^0f^0)p^3q\quad\text{etc.}\\
&-g^0pqq-\tfrac{3}{4}g'ppqq\\
&-(h^0-\tfrac{1}{3}f^0f^0)pq^3
\end{aligned}
\]

It seems worth expanding the area of triangle \(ABD\) in a series as well. For this expansion we use the following condition, which is easily derived from fairly evident geometric considerations and in which \(S\) denotes the area sought:
\[
\frac{r\sin\psi}{n}\mathord{.}\frac{\mathrm{d}S}{\mathrm{d}p}+r\cos\psi\mathord{.}\frac{\mathrm{d}S}{\mathrm{d}q}=\frac{r\sin\psi}{n}\mathord{.}\int n\,\mathrm{d}q
\]
with the integration starting at \(q=0\). We then obtain by the method of undetermined coefficients
\[
\begin{aligned}
\relax[7]\quad S={}&\tfrac{1}{2}pq-\tfrac{1}{12}f^0p^3q-\tfrac{1}{20}f'p^4q-(\tfrac{1}{30}f''-\tfrac{1}{60}f^0f^0)p^5q\quad\text{etc.}\\
&-\tfrac{1}{12}f^0pq^3-\tfrac{3}{40}g^0p^3qq-\tfrac{1}{20}g'p^4qq\\
&-\tfrac{7}{120}f'ppq^3-(\tfrac{1}{15}h^0+\tfrac{2}{45}f''+\tfrac{1}{60}f^0f^0)p^3q^3\\
&-\tfrac{1}{10}g^0pq^4-\tfrac{3}{40}g'ppq^4\\
&-(\tfrac{1}{10}h^0-\tfrac{1}{30}f^0f^0)pq^5
\end{aligned}
\]
\Article{25}
The formulas of the preceding article concern a right triangle formed by shortest lines. We now proceed to the general case. Let \(C\) be another point on the same shortest line
